\documentclass[11pt,a4paper]{article}

% --- Encoding ---
\usepackage[utf8]{inputenc}
\usepackage[T1]{fontenc}
\usepackage[spanish,es-noshorthands]{babel}
\usepackage{lmodern}

% --- Layout ---
\usepackage[margin=2.5cm]{geometry}
\usepackage{parskip}
\setlength{\parindent}{0pt}
\setlength{\parskip}{0.5em plus 0.2em minus 0.1em}

% --- Tables and graphics ---
\usepackage{booktabs}
\usepackage{array}
\usepackage{enumitem}
\usepackage{xcolor}
\usepackage{graphicx}
\usepackage{tikz}
\usetikzlibrary{positioning}
\usepackage{caption}
\usepackage{subcaption}

% --- Code listings ---
\usepackage{listings}
\definecolor{codegray}{rgb}{0.96,0.96,0.96}
\definecolor{codeborder}{rgb}{0.85,0.85,0.85}
\definecolor{codecomment}{rgb}{0.30,0.55,0.30}
\definecolor{codekeyword}{rgb}{0.10,0.30,0.70}
\definecolor{codestring}{rgb}{0.65,0.15,0.10}

\lstdefinestyle{shellstyle}{
    language=bash,
    backgroundcolor=\color{codegray},
    basicstyle=\ttfamily\small,
    breaklines=true,
    frame=single,
    rulecolor=\color{codeborder},
    showstringspaces=false,
    keywordstyle=\color{codekeyword}\bfseries,
    commentstyle=\color{codecomment}\itshape,
    stringstyle=\color{codestring},
    aboveskip=1em,
    belowskip=1em,
    columns=fullflexible,
}

\lstdefinestyle{cstyle}{
    language=C,
    backgroundcolor=\color{codegray},
    basicstyle=\ttfamily\footnotesize,
    breaklines=true,
    frame=single,
    rulecolor=\color{codeborder},
    showstringspaces=false,
    keywordstyle=\color{codekeyword}\bfseries,
    commentstyle=\color{codecomment}\itshape,
    stringstyle=\color{codestring},
    numbers=left,
    numberstyle=\tiny\color{gray},
    numbersep=8pt,
    aboveskip=1em,
    belowskip=1em,
}

\lstset{style=shellstyle}

% --- Hyperlinks ---
\usepackage[colorlinks=true,
            linkcolor=blue!60!black,
            urlcolor=blue!50!black,
            citecolor=blue!60!black]{hyperref}

% --- Title block ---
\title{
    \vspace{-2.5em}
    \textbf{Práctico 1} \\[0.3em]
    \large Comandos, redirección, pipes, señales y llamadas al sistema
}
\author{
    Tomás Rodeghiero \\
    \small Sistemas Operativos --- Universidad Nacional de Río Cuarto
}
\date{2026}

\begin{document}

\maketitle

\begin{abstract}
\noindent
Este documento presenta la resolución integral del Práctico 1 de la
materia Sistemas Operativos (UNRC, 2026). Cada ejercicio se justifica con
la teoría de los capítulos 1, 2 y 3 de las \emph{Notas del curso}
(Marcelo Arroyo) y del capítulo 1 de \emph{Modern Operating Systems}
(Tannenbaum \& Bos). El objetivo es transparentar el modelo
\texttt{fork()}/\texttt{exec()}/\texttt{wait()} junto con los mecanismos de
redirección, pipes, señales y llamadas al sistema que provee la API POSIX.
Todo el código fue compilado con \texttt{gcc -Wall -Wextra} y verificado en
una distribución tipo UNIX.
\end{abstract}

\tableofcontents
\bigskip

% =====================================================================
\section{Introducción}
% =====================================================================

Un \textbf{sistema operativo (SO)} es la capa de software que abstrae el
hardware, administra los recursos disponibles (CPU, memoria, dispositivos
de E/S) y ofrece servicios a las aplicaciones mediante \emph{llamadas al
sistema}. Tannenbaum y Bos, en el capítulo 1 de \emph{Modern Operating
Systems}, lo presentan a la vez como \emph{máquina extendida}---que oculta
la fealdad del hardware tras abstracciones uniformes---y como
\emph{administrador de recursos}, que multiplexa CPU y memoria en tiempo y
en espacio entre múltiples programas.

Las Notas del curso enfatizan tres elementos transversales a este
práctico:

\begin{itemize}[leftmargin=*]
    \item \textbf{Modo kernel y modo usuario.} La CPU ejecuta el código
        del SO con privilegios totales (modo supervisor) y los procesos de
        usuario en modo restringido. El cambio se realiza vía
        \emph{traps} (syscalls) o interrupciones.
    \item \textbf{Proceso.} Instancia de un programa en ejecución, con
        espacio de direcciones propio (código, datos globales, pila),
        descriptor en el kernel (PCB) y archivos abiertos heredados.
    \item \textbf{Multitarea / multiprogramación.} El SO mantiene varios
        procesos en memoria y multiplexa la CPU mediante \emph{context
        switches} (pre-emptive en sistemas de tiempo compartido).
\end{itemize}

Sobre estas bases, el shell se entiende como un intérprete REPL que
compone procesos mediante operadores de redirección (\texttt{<},
\texttt{>}, \texttt{>>}), secuencia (\texttt{;}), condicionales
(\texttt{\&\&}, \texttt{||}), pipes (\texttt{|}) y ejecución concurrente
(\texttt{\&}). Cada operador encuentra una traducción directa en
combinaciones de las syscalls \texttt{fork()}, \texttt{exec()},
\texttt{wait()}, \texttt{pipe()}, \texttt{dup2()} y
\texttt{open/close/read/write}.

% =====================================================================
\section{Ejercicio 1 --- Comandos del shell}
% =====================================================================

\subsection*{1.a Procesos en ejecución y PID del primer proceso}

\begin{lstlisting}
ps -e -o pid,ppid,state,comm | head -n 20
\end{lstlisting}

\texttt{ps -e} lista todos los procesos del sistema. La opción
\texttt{-o} permite seleccionar las columnas relevantes (PID, PPID,
estado y nombre). El proceso con \textbf{PID 1} es el primer proceso
lanzado por el kernel apenas concluye la inicialización del SO; en Linux
suele ser \texttt{systemd} (o \texttt{init} en variantes clásicas), en
macOS \texttt{launchd}. Las Notas del curso lo describen como ``el primer
proceso de usuario preconfigurado'' que se lanza al final de la rutina de
arranque.

\subsection*{1.b Usuarios conectados y usuario actual}

\begin{lstlisting}
who
who am i
whoami
\end{lstlisting}

\texttt{who} lista todas las sesiones abiertas, \texttt{who am i} muestra
únicamente la asociada a la terminal actual y \texttt{whoami} imprime el
\emph{UID efectivo} del proceso shell. La diferencia ilustra la
naturaleza multiusuario de UNIX: el SO mantiene una tabla de sesiones
(\texttt{utmp}/\texttt{wtmp}) que asocia terminales con usuarios
autenticados.

\subsection*{1.c Directorio corriente}

\begin{lstlisting}
pwd
echo "$PWD"
\end{lstlisting}

\texttt{pwd} consulta el directorio asociado al proceso shell, mientras
que \texttt{\$PWD} es una variable de ambiente que el shell mantiene
sincronizada. Ambos comandos cumplen el requisito del enunciado de usar
\emph{al menos dos formas distintas}.

\subsection*{1.d Directorio home del usuario}

\begin{lstlisting}
echo "$HOME"
cd ~
pwd
\end{lstlisting}

\texttt{HOME} es una variable de ambiente fijada en el momento del
\emph{login} a partir de \texttt{/etc/passwd}. El símbolo \texttt{\~{}}
es un atajo que el shell expande a \texttt{\$HOME} antes de ejecutar el
comando.

\subsection*{1.e Metadatos de archivos}

\begin{lstlisting}
ls -la .
ls -la /
stat README.md /
\end{lstlisting}

\texttt{ls -la} muestra entradas ocultas y metadatos básicos (permisos,
dueño, grupo, tamaño, fecha). \texttt{stat} agrega información de bajo
nivel: inodo, cantidad de hard links, timestamps y dispositivo
contenedor. Recordemos que en UNIX un archivo es una secuencia de bytes
identificada por un inodo y la entrada de directorio asocia un nombre con
ese inodo.

\subsection*{1.f Listar contenido sin editor}

\begin{lstlisting}
cat README.md
more README.md
less README.md
\end{lstlisting}

\texttt{cat} vuelca todo el archivo a \texttt{stdout} mediante
\texttt{read()} y \texttt{write()}. \texttt{more} y \texttt{less} agregan
una capa de paginación; \texttt{less} es más flexible al permitir
desplazamiento bidireccional.

\subsection*{1.g Variables de ambiente}

\begin{lstlisting}
env | sort
\end{lstlisting}

Las variables de ambiente forman parte del \emph{contexto} del proceso y
se pasan a los hijos mediante el tercer parámetro de \texttt{main}
(\texttt{char *envv[]}). Las más relevantes son \texttt{USER},
\texttt{HOME}, \texttt{PWD}, \texttt{SHELL}, \texttt{PATH} y
\texttt{LANG}.

\subsection*{1.h Directorios donde el shell busca comandos}

\begin{lstlisting}
echo "$PATH" | tr ':' '\n'
\end{lstlisting}

\texttt{PATH} es una lista de directorios separados por dos puntos. Al
recibir un comando sin barras, el shell recorre esos directorios en
orden y ejecuta el primer archivo regular con permiso de ejecución que
coincida. Por seguridad, el directorio actual \texttt{.} no suele estar
incluido.

\subsection*{1.i Ubicación de un comando con \texttt{which}}

\begin{lstlisting}
which cat
\end{lstlisting}

\texttt{which} reporta el \emph{full path} del primer ejecutable que el
shell encontrará para ese nombre según el \texttt{PATH} vigente.

\subsection*{1.j Crear un archivo con al menos tres comandos}

\begin{lstlisting}
touch archivo.txt
echo "Primera linea" > archivo.txt
date "+%Y-%m-%d %H:%M:%S" >> archivo.txt
echo "Tercera linea" >> archivo.txt
cat archivo.txt
\end{lstlisting}

\texttt{touch} crea (o actualiza) el archivo. El operador \texttt{>}
trunca y reescribe; \texttt{>>} abre en modo \emph{append}. Internamente
el shell, antes de \texttt{exec}, abre el archivo y reasigna el
descriptor 1 con \texttt{dup2()}.

\subsection*{1.k Crear un enlace simbólico}

\begin{lstlisting}
ln -s archivo.txt archivo_link.txt
ls -li archivo.txt archivo_link.txt
\end{lstlisting}

Un enlace simbólico es un archivo cuyo contenido es la \emph{ruta} del
destino. Tiene su propio inodo y \texttt{ls -li} muestra la flecha
\texttt{->} hacia el archivo real.

\subsection*{1.l Eliminar el enlace y comparar con hard link}

\begin{lstlisting}
rm archivo_link.txt
ln archivo.txt archivo_hard.txt
ls -li archivo.txt archivo_hard.txt
\end{lstlisting}

Eliminar el symlink \emph{no} afecta al original: solo se borra la
entrada de directorio que contenía la ruta. Un \emph{hard link}, en
cambio, crea otra entrada que comparte el mismo inodo: el archivo solo
desaparece del disco cuando el contador \texttt{st\_nlink} llega a cero.

\begin{table}[h]
\centering
\small
\begin{tabular}{p{4.5cm} p{5cm} p{5cm}}
\toprule
\textbf{Característica} & \textbf{Symbolic link} & \textbf{Hard link} \\
\midrule
Tipo de referencia & Por ruta (texto) & Por inodo \\
Inodo propio & Sí & No (comparte) \\
Sobrevive al rename del original & Solo si la ruta sigue válida & Sí \\
Sobrevive al borrado del original & No (queda \emph{roto}) & Sí \\
Cruza filesystems & Sí & No \\
Permitido sobre directorios & Sí & Generalmente no \\
\bottomrule
\end{tabular}
\caption{Comparación entre \emph{symbolic links} y \emph{hard links}.}
\end{table}

\subsection*{1.m Programa C que muestra los argumentos}

\begin{lstlisting}[style=cstyle]
#include <stdio.h>

int main(int argc, char *argv[]) {
    printf("argc = %d\n", argc);
    for (int i = 0; i < argc; i++) {
        printf("argv[%d] = %s\n", i, argv[i]);
    }
    return 0;
}
\end{lstlisting}

Compilación y ejecución:

\begin{lstlisting}
gcc -Wall -Wextra -o hello hello.c
./hello arg1 arg2 arg3
\end{lstlisting}

Salida:

\begin{lstlisting}
argc = 4
argv[0] = ./hello
argv[1] = arg1
argv[2] = arg2
argv[3] = arg3
\end{lstlisting}

\paragraph{¿Por qué \texttt{./hello arg1 ... argn}?}
Cuando el usuario escribe un nombre sin barras, el shell asume que se
trata de un comando externo y lo busca \emph{únicamente} en los
directorios listados en \texttt{PATH}. El directorio actual \texttt{.} no
suele estar incluido para evitar el ataque clásico de un binario
malicioso colocado en \texttt{cwd}. El prefijo \texttt{./} indica una
ruta relativa explícita y evita la búsqueda en \texttt{PATH}.

% =====================================================================
\section{Ejercicio 2 --- Redirección, secuencia y filtrado}
% =====================================================================

Los archivos \texttt{correos1.txt} y \texttt{correos2.txt} contienen
direcciones de correo (una por línea) con duplicados intencionales.

\subsection*{2.a Lista ordenada y sin duplicados con archivos temporales}

\begin{lstlisting}
cat correos1.txt correos2.txt > /tmp/correos_todos.txt
sort /tmp/correos_todos.txt > /tmp/correos_ordenados.txt
uniq /tmp/correos_ordenados.txt > correos_ordenados_sin_duplicados.txt
\end{lstlisting}

\texttt{cat} concatena ambos archivos; \texttt{sort} los ordena;
\texttt{uniq} elimina duplicados \emph{consecutivos}. La combinación
\texttt{sort | uniq} elimina globalmente cualquier repetición. Los
temporales se ubican en \texttt{/tmp} porque el SO los purga
periódicamente.

\subsection*{2.b En una sola línea con operador secuencial \texttt{;}}

\begin{lstlisting}
cat correos1.txt correos2.txt > /tmp/correos_todos.txt ; sort /tmp/correos_todos.txt > /tmp/correos_ordenados.txt ; uniq /tmp/correos_ordenados.txt > correos_ordenados_sin_duplicados.txt
\end{lstlisting}

El operador \texttt{;} ejecuta los comandos en secuencia
\emph{independientemente} del éxito o fracaso del anterior. Si quisiéramos
abortar al primer error usaríamos \texttt{\&\&}, pero el enunciado pide
explícitamente operador secuencial.

\subsection*{2.c Filtrar y contar direcciones de \texttt{google.com}}

\begin{lstlisting}
grep -E '@google\.com$' correos_ordenados_sin_duplicados.txt | wc -l
\end{lstlisting}

\begin{itemize}[leftmargin=*]
    \item \texttt{grep -E} interpreta el patrón como expresión regular
        extendida.
    \item \texttt{@google\textbackslash.com\$} fuerza coincidencia
        literal de \texttt{@google}, un punto literal (no un comodín),
        \texttt{com} y el ancla de fin de línea \texttt{\$}.
    \item El pipe envía la salida de \texttt{grep} a \texttt{wc -l}, que
        cuenta líneas.
\end{itemize}

Con los archivos de ejemplo provistos el resultado esperado es \texttt{3}.

% =====================================================================
\section{Ejercicio 3 --- Pipes, sin temporales}
% =====================================================================

\begin{lstlisting}
cat correos1.txt correos2.txt | sort | uniq > correos_ordenados_sin_duplicados.txt
\end{lstlisting}

Las Notas del curso describen al \emph{pipe} como un buffer acotado en el
kernel que implementa el patrón productor/consumidor. Cuando el shell ve
\texttt{cmd1 | cmd2}:

\begin{enumerate}[leftmargin=*]
    \item Crea el pipe con \texttt{pipe(p)}.
    \item Hace \texttt{fork()} para lanzar ambos procesos en paralelo.
    \item Redirige \texttt{stdout} de \texttt{cmd1} al extremo de
        escritura (\texttt{p[1]}) y \texttt{stdin} de \texttt{cmd2} al
        extremo de lectura (\texttt{p[0]}) usando \texttt{dup2()}.
    \item Cierra los descriptores que cada proceso no usa: si quedaran
        abiertos, el lector nunca vería EOF.
\end{enumerate}

\begin{center}
\begin{tikzpicture}[node distance=1.4cm, every node/.style={font=\ttfamily\small}]
\node[draw, rounded corners] (cat) {cat correos*.txt};
\node[draw, rounded corners, right=of cat] (sort) {sort};
\node[draw, rounded corners, right=of sort] (uniq) {uniq};
\node[draw, rounded corners, right=of uniq] (out) {> archivo.txt};
\draw[->, thick] (cat) -- node[above]{|} (sort);
\draw[->, thick] (sort) -- node[above]{|} (uniq);
\draw[->, thick] (uniq) -- (out);
\end{tikzpicture}
\end{center}

A diferencia del ejercicio 2, los datos viajan por buffers en memoria sin
crear archivos en el disco, los tres procesos ejecutan concurrentemente y
no hay limpieza posterior.

% =====================================================================
\section{Ejercicio 4 --- Shell scripts y permisos}
% =====================================================================

\subsection*{4.1 Probar el script con \texttt{sh}}

Archivo \texttt{procesar\_correos\_sin\_interprete.sh} (sin shebang):

\begin{lstlisting}[style=shellstyle]
cat correos1.txt correos2.txt | sort | uniq > correos_ordenados_sin_duplicados.txt
echo "Archivo generado: correos_ordenados_sin_duplicados.txt"
\end{lstlisting}

Invocación:

\begin{lstlisting}
sh procesar_correos_sin_interprete.sh
\end{lstlisting}

Como \texttt{sh} se invoca explícitamente, el archivo no necesita
permiso de ejecución; basta con poder leerlo.

\subsection*{4.2 Agregar shebang y permisos}

\begin{lstlisting}[style=shellstyle]
#!/bin/sh
set -eu
cat correos1.txt correos2.txt | sort | uniq > correos_ordenados_sin_duplicados.txt
echo "Archivo generado: correos_ordenados_sin_duplicados.txt"
\end{lstlisting}

\begin{lstlisting}
chmod +x procesar_correos.sh
./procesar_correos.sh
\end{lstlisting}

El shebang \texttt{\#!/bin/sh} le indica al kernel qué intérprete usar al
invocar \texttt{execve()}. La opción \texttt{set -eu} aborta al primer
error y trata como error variables sin definir; no es exigida por el
enunciado pero es buena práctica defensiva.

% =====================================================================
\section{Ejercicio 5 --- Programa C con \texttt{exit status}}
% =====================================================================

\begin{lstlisting}[style=cstyle]
#include <errno.h>
#include <stdio.h>
#include <stdlib.h>

int main(int argc, char *argv[]) {
    char *endptr = NULL;
    long value;

    if (argc != 2) {
        fprintf(stderr, "Uso: %s <0-255>\n", argv[0]);
        return 2;
    }

    errno = 0;
    value = strtol(argv[1], &endptr, 10);

    if (errno != 0 || endptr == argv[1] || *endptr != '\0') {
        fprintf(stderr, "Error: el argumento debe ser un entero.\n");
        return 2;
    }

    if (value < 0 || value > 255) {
        fprintf(stderr, "Error: exit status valido entre 0 y 255.\n");
        return 2;
    }

    return (int)value;
}
\end{lstlisting}

\paragraph{Justificación.}
Las Notas del curso describen \texttt{exit(code)} como la syscall que
finaliza un proceso \emph{limpiamente} y entrega el código al kernel,
quien luego se lo pasa al padre vía \texttt{wait}/\texttt{waitpid}. Por
convención \texttt{0} significa éxito y cualquier valor distinto, error.
El kernel solo expone los \textbf{8 bits inferiores} del exit code en la
variable \texttt{\$?} del shell, lo que justifica la validación de rango
\texttt{[0,255]}.

\subsection*{5.a Ejecutar y consultar el exit status}

\begin{lstlisting}
./myprog 7
echo $?       # imprime 7
\end{lstlisting}

\subsection*{5.b Ejecutar otro comando si finaliza con éxito}

\begin{lstlisting}
./myprog 0 && echo "myprog termino con exito"
\end{lstlisting}

\texttt{\&\&} dispara el segundo comando solo si el primero retornó
\texttt{0}.

\subsection*{5.c Ejecutar otro comando si finaliza con error}

\begin{lstlisting}
./myprog 5 || echo "myprog termino con error"
\end{lstlisting}

\texttt{||} dispara el segundo comando solo si el primero retornó un
valor distinto de cero.

% =====================================================================
\section{Ejercicio 6 --- Versión propia de \texttt{system()}}
% =====================================================================

La libc estándar provee \texttt{int system(const char *cmd)} como una
abstracción sobre el patrón \texttt{fork()}/\texttt{exec()}/\texttt{wait()}.
Una implementación equivalente es:

\begin{lstlisting}[style=cstyle]
#include <errno.h>
#include <unistd.h>
#include <sys/wait.h>

int mysystem(const char *command) {
    pid_t pid;
    int status;

    if (command == NULL) {
        return access("/bin/sh", X_OK) == 0;
    }

    pid = fork();
    if (pid < 0) return -1;

    if (pid == 0) {
        execl("/bin/sh", "sh", "-c", command, (char *)NULL);
        _exit(127); /* exec falló: convención "comando no ejecutable" */
    }

    while (waitpid(pid, &status, 0) < 0) {
        if (errno != EINTR) return -1;
    }
    return status;
}
\end{lstlisting}

\paragraph{Decisiones de diseño relevantes.}
\begin{itemize}[leftmargin=*]
    \item \textbf{\texttt{/bin/sh -c <cmd>}} en lugar de \texttt{execvp}:
        permite que el comando contenga operadores de shell (pipes,
        redirecciones, \texttt{\&\&}, etc.).
    \item \textbf{\texttt{\_exit(127)}} si \texttt{exec} falla: usa
        \texttt{\_exit} y no \texttt{exit} para evitar que las funciones
        registradas con \texttt{atexit} y los buffers de \texttt{stdio}
        corran dos veces.
    \item \textbf{Reintento ante \texttt{EINTR}}: una señal puede
        interrumpir \texttt{waitpid} antes de que el hijo termine.
\end{itemize}

El \texttt{main} de prueba decodifica el \texttt{status} con
\texttt{WIFEXITED}/\texttt{WEXITSTATUS} para reportar el exit code. Esta
mecánica de dos pasos (\texttt{fork} + \texttt{exec}) es exactamente la
que el shell usa para implementar redirecciones: el padre ajusta el
ambiente del hijo (descriptores, variables) antes de invocar al nuevo
programa.

% =====================================================================
\section{Ejercicio 7 --- Comunicación padre/hijo por \texttt{pipe()}}
% =====================================================================

\paragraph{Idea.}
\texttt{pipe(p)} retorna dos descriptores: \texttt{p[0]} (lectura) y
\texttt{p[1]} (escritura). Como \texttt{fork()} duplica los descriptores,
ambos procesos los heredan; cada uno debe cerrar el extremo que no usa.
Si el padre olvida cerrar el extremo de escritura, el hijo nunca verá
EOF y se bloqueará indefinidamente.

\begin{lstlisting}[style=cstyle]
int p[2];
pipe(p);
pid_t pid = fork();

if (pid == 0) {
    /* Hijo: lee del pipe, imprime en stdout. */
    close(p[1]);
    char buffer[256];
    ssize_t n;
    printf("Hijo recibio: "); fflush(stdout);
    while ((n = read(p[0], buffer, sizeof(buffer))) > 0)
        write(STDOUT_FILENO, buffer, n);
    close(p[0]);
    _exit(0);
}

/* Padre: escribe el mensaje y cierra. */
close(p[0]);
write(p[1], msg, strlen(msg));
close(p[1]);   /* dispara EOF en el hijo */
waitpid(pid, NULL, 0);
\end{lstlisting}

Compilación y ejecución:

\begin{lstlisting}
gcc -Wall -Wextra -o pipe_padre_hijo pipe_padre_hijo.c
./pipe_padre_hijo "Hola desde el padre"
\end{lstlisting}

Salida:

\begin{lstlisting}
Hijo recibio: Hola desde el padre
\end{lstlisting}

\textbf{Limitación:} un pipe sin nombre solo sirve entre procesos
emparentados, porque los descriptores solo se heredan vía \texttt{fork}.
Esa limitación se resuelve en el ejercicio siguiente.

% =====================================================================
\section{Ejercicio 8 --- Comunicación por FIFO}
% =====================================================================

Un \textbf{FIFO} (\emph{named pipe}) es un pipe con un nombre en el
sistema de archivos. Aparece en \texttt{ls -l} con la marca \texttt{p} en
los permisos. Se crea con la syscall \texttt{mkfifo(path, mode)} (o el
comando equivalente) y se abre con la API estándar de archivos. Los datos
\emph{no se almacenan en disco}: solo el \emph{nombre} vive en el
filesystem; el buffer real está en el kernel, igual que un pipe común.

\begin{lstlisting}[style=cstyle]
#define FIFO_PATH "/tmp/practico1_fifo_pipe"

mkfifo(FIFO_PATH, 0600);

if (fork() == 0) {
    int fd = open(FIFO_PATH, O_RDONLY);
    char buf[256]; ssize_t n;
    printf("Hijo recibio: ");
    while ((n = read(fd, buf, sizeof(buf))) > 0)
        write(STDOUT_FILENO, buf, n);
    close(fd); _exit(0);
}

int fd = open(FIFO_PATH, O_WRONLY);
write(fd, msg, strlen(msg));
close(fd);
wait(NULL);
unlink(FIFO_PATH);
\end{lstlisting}

\paragraph{Sincronización implícita.}
\texttt{open(fifo, O\_RDONLY)} se bloquea hasta que algún proceso abra
el otro extremo (\texttt{O\_WRONLY}), y viceversa. Esta propiedad permite
coordinar procesos que se conocen únicamente por la ruta del FIFO.

\paragraph{Limpieza.}
El FIFO sobrevive al cierre de ambos extremos; hay que removerlo
explícitamente con \texttt{unlink()}.

% =====================================================================
\section{Ejercicio 9 --- Mini shell \texttt{minish}}
% =====================================================================

El programa \texttt{minish.c} implementa un mini shell que lee comandos
en bucle hasta recibir \texttt{Ctrl-D} (EOF en \texttt{stdin}) y reconoce
los siguientes operadores binarios:

\begin{table}[h]
\centering
\small
\begin{tabular}{l l}
\toprule
\textbf{Operador} & \textbf{Mecanismo} \\
\midrule
\texttt{cmd1 ; cmd2}     & Ejecutar \texttt{cmd1}, luego \texttt{cmd2}. \\
\texttt{cmd1 \&\& cmd2}  & Ejecutar \texttt{cmd2} solo si \texttt{cmd1} retornó \texttt{0}. \\
\texttt{cmd1 || cmd2}    & Ejecutar \texttt{cmd2} solo si \texttt{cmd1} retornó \texttt{!= 0}. \\
\texttt{cmd1 \& cmd2}    & \texttt{cmd1} en background, \texttt{cmd2} en foreground. \\
\texttt{cmd1 | cmd2}     & \texttt{pipe(p)} + dos \texttt{fork}/\texttt{exec} + \texttt{dup2}. \\
\texttt{cmd < archivo}   & \texttt{open(archivo, O\_RDONLY)} + \texttt{dup2(fd, 0)}. \\
\texttt{cmd > archivo}   & \texttt{open(...,O\_WRONLY|O\_CREAT|O\_TRUNC)} + \texttt{dup2(fd, 1)}. \\
\texttt{cmd >> archivo}  & \texttt{open(...,O\_WRONLY|O\_CREAT|O\_APPEND)} + \texttt{dup2(fd, 1)}. \\
\bottomrule
\end{tabular}
\caption{Operadores reconocidos por \texttt{minish} y su mecanismo POSIX.}
\end{table}

\paragraph{Estructura.}
\begin{enumerate}[leftmargin=*]
    \item \textbf{Loop principal} (\texttt{getline()}). Si retorna $<0$
        (EOF) el shell sale.
    \item \textbf{Tokenización} por espacios y tabulaciones; los tokens
        se almacenan en un arreglo de punteros al buffer original
        modificado in-situ con \texttt{\textbackslash 0}.
    \item \textbf{Despacho recursivo} (\texttt{dispatch}). Busca
        operadores en orden de prioridad baja a alta y descompone el
        comando en \texttt{(left, op, right)} llamando recursivamente a
        cada lado.
    \item \textbf{Comando simple}: \texttt{fork} + (opcional
        \texttt{dup2}) + \texttt{execvp}; el padre llama a
        \texttt{waitpid} salvo en background.
    \item \textbf{Pipe}: dos \texttt{fork} con un \texttt{pipe} previo;
        el padre cierra ambos extremos para que los hijos vean EOF.
\end{enumerate}

\paragraph{Salida esperada.}
\begin{lstlisting}
$ ./minish
minish$ echo hola ; echo mundo
hola
mundo
minish$ sleep 1 & echo "sigo en foreground"
[bg] pid=12345
sigo en foreground
minish$ true && echo "se ejecuto AND"
se ejecuto AND
minish$ false || echo "se ejecuto OR"
se ejecuto OR
minish$ echo hola | tr a-z A-Z
HOLA
minish$ ^D
\end{lstlisting}

\paragraph{Limitaciones.}
El parsing es deliberadamente simple: no soporta comillas, expansión de
variables ni globbing. El propósito didáctico es ilustrar el modelo
\texttt{fork/exec/wait/pipe/dup2}, no reimplementar \texttt{bash}. La
implementación completa se encuentra en \texttt{resolucion-9/minish.c}.

% =====================================================================
\section{Ejercicio 10 --- Captura de \texttt{SIGINT}}
% =====================================================================

\paragraph{Marco teórico.}
Las Notas del curso describen las señales como un mecanismo básico de
comunicación entre procesos para notificar eventos. \texttt{SIGINT}
(interrupción) es la señal por defecto que el shell envía al proceso de
\emph{foreground} cuando el usuario presiona \texttt{Ctrl-C}. Sin
manejador, el kernel termina el proceso. Este ejercicio invierte ese
comportamiento preguntando antes de finalizar.

\begin{lstlisting}[style=cstyle]
#include <ctype.h>
#include <errno.h>
#include <signal.h>
#include <stdio.h>
#include <string.h>
#include <unistd.h>

static void sigint_handler(int sig) {
    const char prompt[] =
        "\nSIGINT recibida. Desea finalizar el proceso? [s/N]: ";
    char answer[64];
    ssize_t nread;
    (void)sig;

    write(STDOUT_FILENO, prompt, sizeof(prompt) - 1);
    do {
        nread = read(STDIN_FILENO, answer, sizeof(answer));
    } while (nread < 0 && errno == EINTR);

    for (ssize_t i = 0; i < nread; i++) {
        unsigned char c = (unsigned char)answer[i];
        if (!isspace(c)) {
            if (c == 's' || c == 'S' || c == 'y' || c == 'Y')
                _exit(0);
            _exit(-1);
        }
    }
    _exit(-1);
}

int main(void) {
    struct sigaction sa = {0};
    sa.sa_handler = sigint_handler;
    sigemptyset(&sa.sa_mask);
    sa.sa_flags = SA_RESTART;
    sigaction(SIGINT, &sa, NULL);

    printf("Proceso en ejecucion. PID=%d\n", (int)getpid());
    for (;;) pause();
}
\end{lstlisting}

\paragraph{Funciones async-signal-safe.}
Las funciones de \texttt{<stdio.h>} no son seguras para invocar dentro de
un handler porque usan buffers y mutexes internos. POSIX define un
subconjunto reducido de funciones \emph{async-signal-safe} (\texttt{read},
\texttt{write}, \texttt{\_exit}, \texttt{sigaction}, etc.) y este código
las usa exclusivamente.

\paragraph{Detalle del exit \texttt{-1}.}
El kernel solo expone los 8 bits inferiores en \texttt{\$?}, así que
\texttt{\_exit(-1)} se observa como \texttt{255}. El programa retorna
lógicamente \texttt{-1}; la representación en el shell es un detalle de
codificación.

% =====================================================================
\section{Ejercicio 11 --- Mínimo global con \texttt{p} hijos y un pipe}
% =====================================================================

\paragraph{Esquema.}
El padre crea \texttt{p} procesos hijos. Cada hijo recibe un bloque de
\texttt{N/p} elementos (la asignación se hace por simple aritmética
\texttt{[i*chunk, (i+1)*chunk)}), calcula su \emph{mínimo local} y lo
envía al padre por un pipe único compartido. El padre lee \texttt{p}
mensajes, va aplicando \texttt{min}, y reporta el resultado tras esperar
a todos los hijos con \texttt{waitpid}.

\begin{lstlisting}[style=cstyle]
typedef struct {
    int child_id;
    int local_min;
} ChildResult;

/* ... validaciones de p y N ... */
pipe(fd);
for (int i = 0; i < p; i++) {
    pid_t pid = fork();
    if (pid == 0) {
        int start = i * chunk, end = start + chunk;
        int local_min = arr[start];
        for (int k = start + 1; k < end; k++)
            if (arr[k] < local_min) local_min = arr[k];

        ChildResult r = { i, local_min };
        close(fd[0]);
        write_full(fd[1], &r, sizeof r);
        close(fd[1]);
        _exit(0);
    }
    children[i] = pid;
}
close(fd[1]);

int global_min = INT_MAX;
for (int i = 0; i < p; i++) {
    ChildResult r;
    read_full(fd[0], &r, sizeof r);
    if (r.local_min < global_min) global_min = r.local_min;
}
close(fd[0]);
for (int i = 0; i < p; i++) waitpid(children[i], NULL, 0);
\end{lstlisting}

\paragraph{¿Por qué un solo pipe basta?}
POSIX garantiza que escrituras de hasta \texttt{PIPE\_BUF} bytes son
\textbf{atómicas} en un pipe: dos escritores nunca mezclan sus mensajes.
Como \texttt{sizeof(ChildResult)} es muy pequeño ($<\!\!16$ bytes), el
padre puede leer mensajes completos en cualquier orden sin corrupción.

\paragraph{Ejemplo.}
\begin{lstlisting}
./minimo_global_pipe 3 8 4 7 10 -2 5
Hijo 0 -> minimo local = 4
Hijo 1 -> minimo local = 7
Hijo 2 -> minimo local = -2
Minimo global = -2
\end{lstlisting}

% =====================================================================
\section{Ejercicio 12 --- Lectura con \texttt{timeout} usando \texttt{alarm}}
% =====================================================================

\paragraph{Marco teórico.}
Las Notas del curso muestran (Listado 6) cómo implementar timeouts: el
kernel envía la señal \texttt{SIGALRM} cuando vence el tiempo programado
con \texttt{alarm(N)}. El handler asociado puede marcar una bandera y, al
volver de la syscall bloqueante, el programa decide si retornó por dato
disponible o por interrupción.

\begin{lstlisting}[style=cstyle]
static volatile sig_atomic_t timed_out = 0;

static void alarm_handler(int sig) {
    (void)sig;
    timed_out = 1;
}

int main(int argc, char *argv[]) {
    int N = atoi(argv[1]);
    char buf[256];

    struct sigaction sa = {0};
    sa.sa_handler = alarm_handler;
    sigemptyset(&sa.sa_mask);
    sa.sa_flags = 0;          /* sin SA_RESTART */
    sigaction(SIGALRM, &sa, NULL);

    alarm(N);
    ssize_t n = read(STDIN_FILENO, buf, sizeof(buf) - 1);
    alarm(0);

    if (n > 0)  { buf[n] = 0; printf("Dato: %s", buf); return 0; }
    if (n == 0) { puts("EOF"); return 0; }
    if (errno == EINTR && timed_out) {
        printf("Timeout: sin datos en %d s\n", N);
        return 124;
    }
    perror("read");
    return 1;
}
\end{lstlisting}

\paragraph{Por qué sin \texttt{SA\_RESTART}.}
Si el handler se instalara con \texttt{SA\_RESTART}, el kernel reanudaría
automáticamente la syscall \texttt{read} después de la señal y nunca
veríamos \texttt{EINTR}. El timeout sería invisible. Sin
\texttt{SA\_RESTART}, \texttt{read} retorna \texttt{-1} con
\texttt{errno = EINTR} y podemos detectar la condición consultando la
bandera \texttt{timed\_out}.

\paragraph{Por qué \texttt{volatile sig\_atomic\_t}.}
\texttt{volatile} impide que el compilador optimice la lectura asumiendo
que la variable no cambia entre dos accesos. \texttt{sig\_atomic\_t} es
el tipo más amplio que el sistema garantiza poder leer y escribir
atómicamente desde un handler.

\paragraph{Por qué \texttt{124} como exit status.}
\texttt{coreutils} usa \texttt{124} como convención para indicar
finalización por timeout, igual que el comando \texttt{timeout(1)}. No es
mágico, pero facilita scripts que distinguen este caso de otros errores.

\paragraph{Pruebas.}
\begin{lstlisting}
(sleep 1; echo "hola") | ./lee_con_timeout 5
Dato recibido: hola

(sleep 3; echo "tarde") | ./lee_con_timeout 1
Timeout: no se recibieron datos en 1 segundos.
echo $?  # imprime 124
\end{lstlisting}

% =====================================================================
\section*{Cierre}
\addcontentsline{toc}{section}{Cierre}
% =====================================================================

A lo largo del práctico se ejercitó la API POSIX como manifestación
concreta de los conceptos teóricos del SO: el modelo de procesos
(\texttt{fork}/\texttt{exec}/\texttt{wait}), la abstracción de archivos y
descriptores (\texttt{open}/\texttt{read}/\texttt{write}/\texttt{dup2}),
los mecanismos de comunicación (\texttt{pipe}, FIFO) y de notificación
asíncrona (\texttt{signal}, \texttt{alarm}). Los operadores del shell
funcionan exactamente como combinaciones de esas syscalls; el mini shell
del ejercicio 9 hace explícita esa correspondencia. La materia en los
prácticos siguientes profundiza en los temas de planificación de CPU,
gestión de memoria y sistemas de archivos que se asoman aquí en su
forma más cruda.

% =====================================================================
\section*{Referencias}
\addcontentsline{toc}{section}{Referencias}
% =====================================================================

\begin{enumerate}[leftmargin=*]
    \item Marcelo Arroyo. \emph{Notas del curso de Sistemas Operativos},
        capítulos 1, 2 y 3. UNRC, 2026.
        \url{https://marceloarroyo.gitlab.io/cursos/SO/}
    \item Andrew S. Tanenbaum, Herbert Bos. \emph{Modern Operating
        Systems}, capítulo 1, 5\textsuperscript{a} edición. Pearson,
        2024.
    \item POSIX.1-2017 / Single UNIX Specification. The Open Group,
        2018.
    \item Sección 2 de \texttt{man} (manuales de las syscalls): \texttt{man
        2 fork}, \texttt{man 2 execve}, \texttt{man 2 pipe}, \texttt{man
        2 sigaction}, etc.
\end{enumerate}

\end{document}
